\documentclass[10pt,a4paper]{amsart}
\usepackage[margin=19mm]{geometry}
\usepackage{amsmath,amssymb,mathtools}
\usepackage[scr=rsfs,cal=euler]{mathalfa}
\usepackage{xfrac}
\usepackage[unicode,hidelinks,bookmarks=false]{hyperref}
\newcommand{\ie}{i.e.\ }
\newcommand{\cf}{cf.\ }
\newcommand{\resp}{respectively\ }
\newcommand{\Subsection}{\subsection}
\providecommand{\nobreakdash}{\nobreak\hbox{-}}
\setlength{\emergencystretch}{2em}
\multlinegap0pt
\usepackage{bm}
\usepackage{bbm}
\usepackage{dsfont}
\usepackage{xspace}
\usepackage{cancel}
\usepackage{mathtools}
\usepackage{amsthm}
\makeatletter
\@ifundefined{theorem}{\newtheorem{theorem}{Theorem}}{}
\@ifundefined{lemma}{\newtheorem{lemma}[theorem]{Lemma}}{}
\@ifundefined{remark}{\newtheorem{remark}[theorem]{Remark}}{}
\makeatother
\title[Bifurcations of viscous boundary layers in the half space]{Bifurcations of viscous boundary layers in the half space}
\author{Dongfen Bian, Emmanuel Grenier, G\'erard Iooss}
\date{Source: arXiv:2510.06715v1 (CC BY 4.0)}
\begin{document}
\maketitle

\noindent\emph{Source and licence.} Bifurcations of viscous boundary layers in the half space. Version arXiv:2510.06715v1 (CC BY 4.0), distributed under the Creative Commons Attribution 4.0 International licence, as recorded in the arXiv licence metadata for this exact version and shown on the abstract page https://arxiv.org/abs/2510.06715v1. Full rights notice: licenses/arxiv-2510.06715-CC-BY-4.0.txt.\par\smallskip

\noindent\textit{Editorial scope.} Counted unit: the Remark whose source label is \texttt{None} in the article (source0.tex lines 1200--1213, source label \texttt{None}), delivered together with the article's own complete original proof of it (source0.tex lines 1215--1344, one \texttt{proof} environment of 130 lines). No mathematics is written, repaired or re-derived: statement and proof are the article's own text line for line. Required context retained uncounted: bound B (lines 1047--1054); quadratic identity (lines 1061--1066); equ diff basic (lines 1187--1189). Every definition, display and label of those lines is retained unchanged. Omitted from this package and not redistributed: 1--1046, 1055--1060, 1067--1186, 1190--1199, 1214--1214, 1345--4970. Nothing inside a retained excerpt is omitted or altered: every retained body is delivered exactly as the article's lines are written, so no omission receipt of any kind accompanies this package. Citations: the retained excerpts carry no citation command at all, so no bibliography entry of the article is transcribed and no reference is redistributed. Reference adaptation: none was needed and none was made; every reference inside the delivered unit resolves inside it, so no unresolved reference or citation is produced. Printed slips kept literal: no printed word, symbol, hypothesis or number of the counted statement or of its proof was altered. Layout and class: the article's own class is \texttt{article} with packages including amsfonts, amssymb, amsmath, amsthm, graphicx, bm, xcolor; the wrapper is the lane's pinned \texttt{amsart} editorial wrapper at 10pt on A4, which restates the theorem-like environments the retained text needs and adds the article's own macro definitions that the retained text needs (none), with extra wrapper packages (bm, bbm, dsfont, xspace, cancel, mathtools). The delivered numbering is the unit's own. Assets: no retained span contains a figure environment, a graphics-inclusion command, a PDF-page inclusion, a table or a TikZ picture; no image file is copied into this package, the package declares no asset dependency and no image file is redistributed with it. Licence: version arXiv:2510.06715v1 of this article is distributed under the Creative Commons Attribution 4.0 International licence, as recorded in the arXiv licence metadata for this exact version and shown on the abstract page https://arxiv.org/abs/2510.06715v1. A full rights notice is delivered at licenses/arxiv-2510.06715-CC-BY-4.0.txt. Characters: every retained excerpt is pure ASCII, so the delivered engine, XeLaTeX with its default Latin Modern text font, reports no missing character.
\begin{lemma}
The quadratic operator $\boldsymbol{B}$ is bounded from $\mathcal{Z}_{\eta}$ to 
$\mathcal{Y}_{\eta}:$ there exists $C>0$ such that%
\begin{equation}
\|\boldsymbol{B}(u,v)\|_{\mathcal{Y}_{\eta}}\leq C\|u\|_{\mathcal{Z}_{\eta}}\|v\|_{%
\mathcal{Z}_{\eta}}.  \label{bound B}
\end{equation}
\end{lemma}

\begin{eqnarray}
\widetilde{P}B(u,v) &=&\widetilde{P}B(\widetilde{u},\widetilde{v})+B(%
\widetilde{u},v_{0})+B(\widetilde{v},u_{0}),  \notag \\
P_{0}B(u,v) &=&P_{0}B(\widetilde{u},\widetilde{v}).
\label{quadratic identity}
\end{eqnarray}%

\begin{equation}
\frac{dv}{dt}=\boldsymbol{L}_{(\nu )}v+B(v,v),  \label{equ diff basic}
\end{equation}%

\begin{remark}
Note that this Theorem implies that 
\begin{eqnarray*}
\|\widetilde{P}v(t)\|_{\mathring{H}_{\eta}^{2}} &\leq &Me^{-\xi t}\|%
\widetilde{P}v(0)\|_{\mathring{H}_{\eta}^{2}},\text{ }t\geq 0, \\
\|P_{0}v(t)\|_{\mathring{W}^{2,\infty }} &\leq &\frac{M}{1+t}%
\Bigl( \|P_{0}v(0)\|_{\mathring{W}_{\eta}^{2}}+\|\widetilde{P}v(0)\|_{\mathring{H}%
_{\eta}^{2}}^{2} \Bigr),\text{ }t\geq 0.
\end{eqnarray*}%
The slow decrease of the $0$-mode as $t\rightarrow \infty $ is due to the
influence of the essential spectrum of $\boldsymbol{L}_{(\nu )}$ which contains
the full negative real line. Moreover, we may notice that the decay in $y$
at $\infty $ is lost for the $0$-mode for $t>0.$
\end{remark}

\begin{proof}
Using (\ref{quadratic identity}) and (\ref{bound B}) shows that there exists 
$C>0$ such that we have the estimates%
\begin{eqnarray*}
\Bigl\|\widetilde{P}B(\widetilde{v},\widetilde{v})+2B(\widetilde{v},v_{0})\Bigr\|_{%
\mathring{H}_{\eta}^{1}} &\leq &C \Bigl(\|\widetilde{v}\|_{\mathring{H}%
_{\eta}^{2}}^{2}+\|\widetilde{v}\|_{\mathring{H}_{\eta}^{2}}\|v_{0}\|_{\mathring{W}%
^{2,\infty }} \Bigr), \\
\|P_{0}B(\widetilde{v},\widetilde{v})\|_{\mathring{W}_{\eta}^{1,\infty }} &\leq
&C\|\widetilde{v}\|_{\mathring{H}_{\eta}^{2}}^{2}.
\end{eqnarray*}
Now (\ref{equ diff basic}) becomes%
\begin{eqnarray*}
\frac{d\widetilde{v}}{dt} &=&\boldsymbol{L}_{(\nu )}\widetilde{v}+\widetilde{%
P}B(\widetilde{v},\widetilde{v})+2B(\widetilde{v},v_{0}), \\
\frac{dv_{0}}{dt} &=&\boldsymbol{L}_{(\nu )}v_{0}+P_{0}B(\widetilde{v},%
\widetilde{v}),
\end{eqnarray*}%
which leads to the following integral formulation for $t\geq 0$%
\begin{eqnarray}
\widetilde{v}(t) &=&e^{\boldsymbol{L}_{(\nu )}t}\widetilde{P}%
v(0)+\int_{0}^{t}e^{\boldsymbol{L}_{(\nu )}(t-s)} \Bigl[\widetilde{P}B(\widetilde{v%
},\widetilde{v})+2B(\widetilde{v},v_{0}) \Bigr](s)ds,  \label{compvtilde} \\
v_{0}(t) &=&e^{\boldsymbol{L}_{(\nu )}t}P_{0}v(0)+\int_{0}^{t}e^{\boldsymbol{%
L}_{(\nu )}(t-s)}P_{0}B(\widetilde{v},\widetilde{v})(s)ds,  \label{compv0}
\end{eqnarray}%
where we look for $\widetilde{v}+v_{0}\in \mathcal{Z}_{\eta,\xi }.$ 
We can solve the system above by the implicit function theorem with respect
to $(\widetilde{v},v_{0})$ in the neighborhood of $0$ for $v(0)=\widetilde{P}%
v(0)+P_{0}v(0)$ provided 
\begin{eqnarray*}
&&\widetilde{P}v(0)\text{ sufficiently small in }\mathring{H}_{\eta}^{2}, \\
&&P_{0}v(0)\text{ sufficiently small in }\mathring{W}_{\eta}^{2}.
\end{eqnarray*}
Using the bound of the imaginary parts of eigenvalues which are isolated in
the left half complex plane, there exist $\xi_{m}>0$
such that all the eigenvalues $\lambda $ of $\boldsymbol{L}_{(\nu )}$ satisfy%
\begin{equation*}
\sup \Re \lambda =-\xi_{m} <-\xi  <0.
\end{equation*}%
We choose $\xi_1$ such that $- \xi_m <  - \xi_1 < -  \xi$.
Then we have the estimates%
\begin{eqnarray*}
\|e^{\boldsymbol{L}_{(\nu )}t}\|_{\mathcal{L}(\mathring{H}_{\eta}^{2})} &\leq
&Ce^{-\xi _{1}t},\text{ }t\geq 0, \\
\|e^{\boldsymbol{L}_{(\nu )}t}\|_{\mathcal{L}(\mathring{W}_{\eta}^{2,\infty },%
\mathring{W}^{2,\infty })} &\leq &\frac{C}{1+t},\text{ }t\geq 0, \\
\|e^{\boldsymbol{L}_{(\nu )}(t-s)}\|_{\mathcal{L}(\mathring{H}_{\eta}^{1},%
\mathring{H}_{\eta}^{2})} &\leq &C \Bigl[1+\frac{1}{\sqrt{t-s}} \Bigr] e^{-\xi_{1}(t-s)},\text{ }t>s, \\
\|e^{\boldsymbol{L}_{(\nu )}(t-s)}\|_{\mathcal{L}(\mathring{W}_{\eta}^{1,\infty
},\mathring{W}^{2,\infty })} &\leq &\frac{C}{\sqrt{t-s}(1 + \sqrt{t-s})},\text{ }t>s.
\end{eqnarray*}%
We note that%
\begin{eqnarray*}
\|e^{\xi t}e^{\boldsymbol{L}_{(\nu )}t}\widetilde{P}v(0)\|_{\mathring{H}%
_{\eta}^{2}} &\leq &C\|\widetilde{P}v(0)\|_{\mathring{H}_{\eta}^{2}}, \\
\|e^{\boldsymbol{L}_{(\nu )}t}P_{0}v(0)\|_{\mathring{W}^{2,\infty }} &\leq &%
\frac{C}{1+t}\|P_{0}v(0)\|_{\mathring{W}_{\eta}^{2}},
\end{eqnarray*}%
that
\begin{eqnarray*}
&&\Bigl\|e^{\xi t}e^{\boldsymbol{L}_{(\nu )}(t-s)} \Bigl[\widetilde{P}B(\widetilde{v}%
,\widetilde{v})+2B(\widetilde{v},v_{0})\Bigr](s)\Bigr\|_{\mathring{H}_{\eta}^{2}} \\
&\leq &MC \Bigl(1+\frac{1}{\sqrt{t-s}} \Bigr)e^{-(\xi _{1}-\xi )(t-s)}|||%
\widetilde{v}|||_{\eta,\xi }|||\widetilde{v}+v_{0}|||_{\eta,\xi },
\end{eqnarray*}%
and that 
\begin{equation*}
\Bigl\|e^{\boldsymbol{L}_{(\nu )}(t-s)}P_{0}B(\widetilde{v},\widetilde{v})(s) \Bigr\|_{%
\mathring{W}^{2,\infty }}\leq \frac{MC}{\sqrt{t-s} (1 + \sqrt{t-s})}e^{-2\xi s}|||%
\widetilde{v}|||_{\eta,\xi }^{2}.
\end{equation*}%
Hence the right hand side of (\ref{compvtilde}), and (\ref{compv0}) is
quadratic and $C^{1}$ in $(\widetilde{v},v_{0})\in \mathcal{Z}_{\eta,\xi }.$
It results by the implicit function theorem, that, provided $(\widetilde{P}%
v(0),P_{0}v(0))$ small enough in $\mathring{H}_{\eta}^{2}\oplus \mathring{W}%
_{\eta}^{2},$ there exists a unique solution $(\widetilde{v},v_{0})\in \mathcal{%
Z}_{\eta,\xi}$ in a neighborhood of $0.$ Moreover, directly from the system
(\ref{compvtilde}), (\ref{compv0}), we have the estimate%
\begin{align*}
\|e^{\xi  t}\widetilde{v}(t)\|_{\mathring{H}_{\eta}^{2}}
\leq C\|\widetilde{P}%
v(0)\|_{\mathring{H}_{\eta}^{2}}+MC\int_{0}^{t} &(1+\frac{1}{(t-s)^{1/2}}%
)
\\ 
& \times e^{-(\xi _{1}-\xi  )(t-s)}|||\widetilde{v}|||_{\eta,\xi  }|||\widetilde{%
v}+v_{0}|||_{\eta,\xi } ds,
\end{align*}%
hence%
\begin{equation*}
|||\widetilde{v}|||_{\eta,\xi }\leq C\|\widetilde{P}v(0)\|_{\mathring{H}%
_{\eta}^{2}}+K|||\widetilde{v}|||_{\eta,\xi }|||\widetilde{v}%
+v_{0}|||_{\eta,\xi  },
\end{equation*}%
using%
\begin{equation*}
\int_{0}^{t} \Bigl[ 1+\frac{1}{(t-s)^{1/2}} \Bigr] e^{-\delta (t-s)}ds\leq C_{1}(\delta )%
\text{ for }\delta >0.
\end{equation*}%
In the same way, we obtain%
\begin{equation*}
|||v_{0}|||_{\eta,\xi }\leq C\|P_{0}v(0)\|_{\mathring{W}_{\eta}^{2,\infty
}}+K|||\widetilde{v}|||_{\eta,\xi }^{2},
\end{equation*}%
using%
\begin{equation*}
\int_{0}^{t}\frac{1}{\sqrt{t-s}) (1+\sqrt{t-s})}e^{-\delta s}ds\leq \frac{C_{2}(\delta )}{1+t},\text{ for }\delta >0.
\end{equation*}%
Hence, 
\begin{equation*}
|||\widetilde{v}|||_{\eta,\xi }\leq C\|\widetilde{P}v(0)\|_{\mathring{H}%
_{\eta}^{2}}+CK\|P_{0}v(0)\|_{\mathring{W}_{\eta}^{2,\infty }}|||\widetilde{v}%
|||_{\eta,\xi }+K|||\widetilde{v}|||_{\eta,\xi }^{2}+K^{2}|||\widetilde{v}%
|||_{\eta,\xi }^{3},
\end{equation*}%
so that for $v(0)$ such that%
\begin{equation*}
CK\|P_{0}v(0)\|_{\mathring{W}_{\eta}^{2,\infty }}+2CK\|\widetilde{P}v(0)\|_{%
\mathring{H}_{\eta}^{2}}+4C^{2}K^{2}\|\widetilde{P}v(0)\|_{\mathring{H}%
_{\eta}^{2}}^{2}< \frac{1}{2},
\end{equation*}%
we obtain%
\begin{eqnarray*}
|||\widetilde{v}|||_{\eta,\xi } &\leq &2C\|\widetilde{P}v(0)\|_{\mathring{H}%
_{\eta}^{2}}, \\
|||v_{0}|||_{\eta,\xi } &\leq &C\|P_{0}v(0)\|_{\mathring{W}_{\eta}^{2,\infty
}}+4C^{2}K\|\widetilde{P}v(0)\|_{\mathring{H}_{\eta}^{2}}^{2}
\end{eqnarray*}%
which ends the proof of the nonlinear asymptotic stability.
\end{proof}

\end{document}
